\section{Circuit Complexity --- Tree DAG}
%%% generated by the blueprint pipeline from TCSlib.Complexity.CircuitComplexity.TreeDAG
\subsection*{Circuit Complexity --- Tree DAG: Introduction}

Declarations of the Lean module \texttt{TCSlib.Complexity.CircuitComplexity.TreeDAG}.

\subsection*{Circuit Complexity --- Tree DAG: Results}

\begin{definition}[Value of a circuit vertex]
\lean{BoolCircuit.vertexValue}
\label{BoolCircuit.vertexValue}
\leanok
\uses{BoolCircuit.DAGGate, BoolCircuit.DAGGate.eval, BoolCircuit.runWith}
Let $g_1,\ldots,g_m$ be Boolean gates, let $x$ be a Boolean assignment to the $n$ input
vertices, and let $v$ be a natural-numbered vertex. Evaluate the gates in order after
initializing the vertex values with $x$, and take the value at index $v$; if that index
is outside the resulting list of values, the value is false.
\end{definition}

\begin{definition}[Depth of a circuit vertex]
\lean{BoolCircuit.vertexDepth}
\label{BoolCircuit.vertexDepth}
\leanok
\uses{BoolCircuit.DAGGate, BoolCircuit.DAGGate.depth, BoolCircuit.runWith}
Let $n$ be the number of input vertices, let $g_1,\ldots,g_m$ be Boolean gates, and let
$v$ be a vertex index. Evaluate the gates in order starting with depth $0$ at each input
vertex, and take the depth at index $v$; if that index is outside the resulting list of
depths, return $0$.
\end{definition}

\begin{theorem}[Vertex values after appending gates]
\lean{BoolCircuit.vertexValue_append}
\label{BoolCircuit.vertexValue_append}
\leanok
\uses{BoolCircuit.DAGGate, BoolCircuit.runWith_append, BoolCircuit.runWith_getD_of_lt, BoolCircuit.vertexValue}
\difficulty{3}
Let $n$ be a natural number, let $gs$ and $ext$ be lists of gates, and let $x$ be an
assignment of Boolean values to the $n$ inputs. For any vertex $v<n+|gs|$, its value
after applying $gs$ followed by $ext$ to $x$ equals its value after applying $gs$ to
$x$.
\end{theorem}

\begin{theorem}[Vertex Depth Under Appending Gates]
\lean{BoolCircuit.vertexDepth_append}
\label{BoolCircuit.vertexDepth_append}
\leanok
\uses{BoolCircuit.DAGGate, BoolCircuit.runWith_append, BoolCircuit.runWith_getD_of_lt, BoolCircuit.vertexDepth}
\difficulty{3}
Let $n$ be a natural number, and let $gs$ and $ext$ be finite lists of gates. For every
natural-number vertex index $v<n+|gs|$, the depth assigned to $v$ after processing $gs$
followed by $ext$, starting with $n$ input vertices of depth zero, equals its depth
after processing $gs$ alone.
\end{theorem}

\begin{theorem}[Value of the last vertex]
\lean{BoolCircuit.vertexValue_last}
\label{BoolCircuit.vertexValue_last}
\leanok
\uses{BoolCircuit.DAGGate, BoolCircuit.DAGGate.eval, BoolCircuit.runWith, BoolCircuit.runWith_getD_last, BoolCircuit.vertexValue}
\difficulty{3}
Let $n$ be a natural number, let $gs$ be a list of gates, and let $g$ be a gate. For any
Boolean assignment $x$ to the $n$ inputs, the value at vertex $n+|gs|$ after applying
$gs$ followed by $g$ equals the evaluation of $g$ on the vertex values obtained by
applying $gs$ to $x$.
\end{theorem}

\begin{theorem}[Depth of the last vertex]
\lean{BoolCircuit.vertexDepth_last}
\label{BoolCircuit.vertexDepth_last}
\leanok
\uses{BoolCircuit.DAGGate, BoolCircuit.DAGGate.depth, BoolCircuit.runWith, BoolCircuit.runWith_getD_last, BoolCircuit.vertexDepth}
\difficulty{3}
Let $n$ be a natural number, let $gs$ be a list of directed acyclic graph gates, and let
$g$ be a gate. In the circuit formed by applying $gs$ followed by $g$ to $n$ initial
vertices, the depth of vertex $n+|gs|$ equals the depth of $g$ computed from the vertex
depths after applying $gs$ to those initial vertices.
\end{theorem}

\begin{theorem}[Value of an input vertex]
\lean{BoolCircuit.vertexValue_input}
\label{BoolCircuit.vertexValue_input}
\leanok
\uses{BoolCircuit.DAGGate, BoolCircuit.runWith_getD_of_lt, BoolCircuit.vertexValue}
\difficulty{3}
Let $n$ be a natural number, let $gs$ be a list of gates, let $x$ assign Boolean values
to the $n$ inputs, and let $i$ be one of those inputs. The value at input vertex $i$
after applying $gs$ to $x$ is $x(i)$.
\end{theorem}

\begin{theorem}[Depth of an input vertex]
\lean{BoolCircuit.vertexDepth_input}
\label{BoolCircuit.vertexDepth_input}
\leanok
\uses{BoolCircuit.DAGGate, BoolCircuit.runWith_getD_of_lt, BoolCircuit.vertexDepth}
\difficulty{3}
Let $n$ be a natural number, let $gs$ be a list of directed acyclic graph gates, and let
$i<n$. In the circuit formed by applying the gates in $gs$ to $n$ initial input
vertices, the depth assigned to vertex $i$ is $0$.
\end{theorem}

\begin{theorem}[Deduplication preserves a list-wide test]
\lean{BoolCircuit.all_dedup}
\label{BoolCircuit.all_dedup}
\leanok
\difficulty{3}
For every list $l$ of natural numbers and every Boolean-valued function $p$ on the
natural numbers, applying the test that every entry satisfies $p$ gives the same result
on $l$ and on the list obtained by removing repeated entries from $l$.
\end{theorem}

\begin{theorem}[Disjunction is unchanged by deduplication]
\lean{BoolCircuit.any_dedup}
\label{BoolCircuit.any_dedup}
\leanok
\difficulty{3}
For every list $l$ of natural numbers and every Boolean-valued predicate $p$, the
listwise disjunction of $p$ over the duplicate-free version of $l$ equals its listwise
disjunction over $l$.
\end{theorem}

\begin{theorem}[Bound on a folded maximum]
\lean{BoolCircuit.foldr_max_le}
\label{BoolCircuit.foldr_max_le}
\leanok
\difficulty{4}
Let $l$ be a finite list of nonnegative integers, let $f$ map nonnegative integers to
nonnegative integers, and let $B$ be a nonnegative integer. If $f(a)\le B$ for every $a$
in $l$, then the maximum formed by folding the values $f(a)$ over $l$, starting from
$0$, is at most $B$.
\end{theorem}

\begin{theorem}[Evaluation of an AND node]
\lean{BoolCircuit.TreeCircuit.eval_node_true}
\label{BoolCircuit.TreeCircuit.eval_node_true}
\leanok
\uses{BoolCircuit.TreeCircuit, BoolCircuit.TreeCircuit.eval, BoolCircuit.TreeCircuit.eval_node_true_iff}
\difficulty{3}
Let $n$ be a natural number, let $cs$ be a list of Boolean circuit trees on $n$
variables, and let $x$ be an assignment of Boolean values to those variables. The
circuit consisting of an AND node with children $cs$ evaluates under $x$ to the same
Boolean value as the conjunction of the evaluations of all circuits in $cs$.
\end{theorem}

\begin{theorem}[Evaluation of an OR node]
\lean{BoolCircuit.TreeCircuit.eval_node_false}
\label{BoolCircuit.TreeCircuit.eval_node_false}
\leanok
\uses{BoolCircuit.TreeCircuit, BoolCircuit.TreeCircuit.eval, BoolCircuit.TreeCircuit.eval_node_false_iff}
\difficulty{3}
For every natural number $n$, every list $cs$ of Boolean circuit trees on $n$ variables,
and every assignment $x : \mathrm{Fin}\,n \to \bbf_2$, the value of the OR node with
children $cs$ is equal to the Boolean disjunction of the values of the circuits in $cs$
under $x$.
\end{theorem}

\begin{definition}[Boolean node gate kind]
\lean{BoolCircuit.nodeKind}
\label{BoolCircuit.nodeKind}
\leanok
\uses{BoolCircuit.GateKind}
For a Boolean flag $b$, assign the AND gate kind when $b$ is true and the OR gate kind
when $b$ is false.
\end{definition}

\begin{definition}[Tree circuit compilation]
\lean{BoolCircuit.compileTree}
\label{BoolCircuit.compileTree}
\leanok
\uses{BoolCircuit.DAGGate, BoolCircuit.TreeCircuit, BoolCircuit.nodeKind}
Given a Boolean circuit tree on $n$ variables and an existing list of gates, compile the
tree by processing its children in order. A positive literal returns its variable index
and leaves the gate list unchanged; a negative literal appends a NOT gate reading that
variable and returns the index $n$ plus the resulting gate-list length. At a node,
append an AND or OR gate according to its Boolean label, reading the deduplicated list
of compiled child outputs, and return $n$ plus the final gate-list length as the output
index.
\end{definition}

\begin{definition}[List of tree circuit compilations]
\lean{BoolCircuit.compileTrees}
\label{BoolCircuit.compileTrees}
\leanok
\uses{BoolCircuit.DAGGate, BoolCircuit.TreeCircuit}
Given a list of Boolean circuit trees on $n$ variables and an existing list of gates,
compile the trees in order, passing the gates produced for each tree as the starting
list for the next. Return the final gate list together with the list of output indices,
in the original tree order.
\end{definition}

\begin{definition}[Tree circuit compilation specification]
\lean{BoolCircuit.CompileSpec}
\label{BoolCircuit.CompileSpec}
\leanok
\uses{BoolCircuit.DAGGate, BoolCircuit.DAGGate.WellFormed, BoolCircuit.GatesAcyclic, BoolCircuit.TreeCircuit, BoolCircuit.TreeCircuit.depth, BoolCircuit.TreeCircuit.eval, BoolCircuit.TreeCircuit.maxFanin, BoolCircuit.TreeCircuit.size, BoolCircuit.vertexDepth, BoolCircuit.vertexValue}
Let $n$ be a natural number, let $c$ be a Boolean circuit tree on $n$ variables, let $g$
be a finite list of gates, and let $(g',v)$ consist of a finite list of gates and a
natural-number vertex. The specification requires that $g'$ be obtained by appending a
list of gates to $g$, that $g'$ be acyclic, and that $v<n+|g'|$; for every input $x$,
the value at vertex $v$ in $g'$ must equal the evaluation of $c$ at $x$. It further
requires $|g'|\le |g|+\mathrm{size}(c)$, that the depth at $v$ be at most
$\mathrm{depth}(c)+1$, and that every gate in the appended portion be well formed and
have at most $\max(1,\mathrm{maxFanin}(c))$ arguments.
\end{definition}

\begin{definition}[List compilation specification]
\lean{BoolCircuit.CompileListSpec}
\label{BoolCircuit.CompileListSpec}
\leanok
\uses{BoolCircuit.DAGGate, BoolCircuit.DAGGate.WellFormed, BoolCircuit.GatesAcyclic, BoolCircuit.TreeCircuit, BoolCircuit.TreeCircuit.eval, BoolCircuit.TreeCircuit.maxDepth, BoolCircuit.TreeCircuit.maxFaninL, BoolCircuit.TreeCircuit.sumSize, BoolCircuit.vertexDepth, BoolCircuit.vertexValue}
Let $n$ be a natural number, let $C$ be a list of Boolean circuit trees on $n$
variables, let $G$ be an initial list of gates, and let $(G',V)$ be a pair consisting of
a gate list and a list of output vertices. The specification requires that $G'$ be
obtained by appending gates to $G$ and be acyclic; that $V$ have the same length as $C$,
with every $v\in V$ satisfying $v<n+|G'|$, and that for every input the values at the
vertices of $V$, in order, equal the evaluations of the trees in $C$. It further
requires $|G'|\le |G|+\operatorname{sumSize}(C)$, every output vertex to have depth at
most $\operatorname{maxDepth}(C)+1$, and every appended gate to be well formed and to
have at most $\max(1,\operatorname{maxFanin}(C))$ arguments.
\end{definition}

\begin{theorem}[Dropping an Appended List]
\lean{BoolCircuit.drop_append_of_prefix}
\label{BoolCircuit.drop_append_of_prefix}
\leanok
\uses{BoolCircuit.DAGGate}
\difficulty{1}
For any two lists of directed acyclic graph gates, dropping from their concatenation the
number of entries in the first list leaves the second list.
\end{theorem}

\begin{theorem}[Specification of list compilation]
\lean{BoolCircuit.compileTrees_spec}
\label{BoolCircuit.compileTrees_spec}
\leanok
\uses{BoolCircuit.CompileListSpec, BoolCircuit.CompileSpec, BoolCircuit.GatesAcyclic, BoolCircuit.TreeCircuit, BoolCircuit.TreeCircuit.maxDepth_cons, BoolCircuit.TreeCircuit.maxFaninL_cons, BoolCircuit.TreeCircuit.sumSize, BoolCircuit.TreeCircuit.sumSize_cons, BoolCircuit.TreeCircuit.sumSize_nil, BoolCircuit.compileTree, BoolCircuit.compileTrees, BoolCircuit.vertexDepth_append, BoolCircuit.vertexValue_append}
\difficulty{6}
Let $n$ be a natural number and $cs$ a finite list of Boolean circuit trees on $n$
variables. Suppose that every $c$ in $cs$, when compiled after any acyclic gate list
$gs$, satisfies the single-circuit compilation specification: the compiled list extends
$gs$ and is acyclic; its root index is below $n$ plus the resulting list length; its
root value agrees with $c$ under every variable assignment; its length is at most $|gs|$
plus the size of $c$; its root depth is at most the depth of $c$ plus $1$; and its new
gates have distinct inputs, every NOT gate has one input, and each new gate has at most
$\max(1,F)$ inputs, where $F$ is the maximum fanin of $c$. Then, for every acyclic
initial gate list $gs$, compiling all of $cs$ appends a gate list and produces a list of
root vertices of the same length as $cs$. The resulting gate list is acyclic and extends
$gs$; every root index is below $n$ plus the resulting list length, and under each
variable assignment the root values, in order, equal the evaluations of the circuits in
$cs$. The resulting gate list has length at most $|gs|$ plus the sum of the circuit
sizes in $cs$, and every root has depth at most the maximum circuit depth in $cs$ plus
$1$. Each new gate has distinct inputs, every NOT gate has one input, and each new gate
has at most $\max(1,F)$ inputs, where $F$ is the maximum fanin among circuits in $cs$.
\end{theorem}

\begin{theorem}[Specification of circuit compilation]
\lean{BoolCircuit.compileTree_spec}
\label{BoolCircuit.compileTree_spec}
\leanok
\uses{BoolCircuit.CompileSpec, BoolCircuit.DAGGate, BoolCircuit.DAGGate.depth, BoolCircuit.DAGGate.eval, BoolCircuit.GatesAcyclic, BoolCircuit.GatesAcyclic.snoc, BoolCircuit.Lit.eval, BoolCircuit.TreeCircuit, BoolCircuit.TreeCircuit.depth, BoolCircuit.TreeCircuit.depth_node, BoolCircuit.TreeCircuit.eval, BoolCircuit.TreeCircuit.eval_node_false, BoolCircuit.TreeCircuit.eval_node_true, BoolCircuit.TreeCircuit.ind, BoolCircuit.TreeCircuit.maxDepth, BoolCircuit.TreeCircuit.maxFanin_node, BoolCircuit.TreeCircuit.size, BoolCircuit.TreeCircuit.size_node, BoolCircuit.all_dedup, BoolCircuit.any_dedup, BoolCircuit.compileTree, BoolCircuit.compileTrees, BoolCircuit.compileTrees_spec, BoolCircuit.drop_append_of_prefix, BoolCircuit.foldr_max_le, BoolCircuit.nodeKind, BoolCircuit.runWith, BoolCircuit.vertexDepth, BoolCircuit.vertexDepth_input, BoolCircuit.vertexDepth_last, BoolCircuit.vertexValue, BoolCircuit.vertexValue_input, BoolCircuit.vertexValue_last}
\difficulty{6}
Let $n$ be a natural number, $c$ a Boolean circuit tree on $n$ variables, and $gs$ a
finite list of gates such that every input to its $i$th gate has index less than $n+i$.
Compiling $c$ after $gs$ produces an extended gate list and a root vertex. The resulting
list extends $gs$, is acyclic in the same sense, and the root index is less than $n$
plus the length of the resulting list. For every assignment of Boolean values to the
variables, the value at the root equals the value of $c$. The resulting list has length
at most $|gs|$ plus the size of $c$, and the depth at the root is at most the depth of
$c$ plus $1$. Each newly appended gate has distinct input indices, any NOT gate among
them has exactly one input, and each has at most $\max(1,F)$ inputs, where $F$ is the
maximum fanin of $c$.
\end{theorem}

\begin{definition}[DAG representation of a tree circuit]
\lean{BoolCircuit.TreeCircuit.toDAG}
\label{BoolCircuit.TreeCircuit.toDAG}
\leanok
\uses{BoolCircuit.DAGCircuit, BoolCircuit.GatesAcyclic.nil, BoolCircuit.TreeCircuit, BoolCircuit.compileTree, BoolCircuit.compileTree_spec}
Every Boolean circuit tree $c$ on $n$ variables determines a directed acyclic circuit
whose gates are the gates obtained by compiling $c$ from an empty gate list and whose
output is the resulting output vertex.
\end{definition}

\begin{theorem}[Evaluation preserved by the DAG translation]
\lean{BoolCircuit.TreeCircuit.toDAG_eval}
\label{BoolCircuit.TreeCircuit.toDAG_eval}
\leanok
\uses{BoolCircuit.DAGCircuit.eval, BoolCircuit.GatesAcyclic.nil, BoolCircuit.TreeCircuit, BoolCircuit.TreeCircuit.eval, BoolCircuit.TreeCircuit.toDAG, BoolCircuit.compileTree_spec}
\difficulty{1}
For every Boolean circuit $c$ on $n$ variables and every Boolean assignment $x$ to those
variables, evaluating the directed acyclic graph translation of $c$ at $x$ gives the
same value as evaluating $c$ at $x$.
\end{theorem}

\begin{theorem}[Size bound for compiled circuits]
\lean{BoolCircuit.TreeCircuit.toDAG_size_le}
\label{BoolCircuit.TreeCircuit.toDAG_size_le}
\leanok
\uses{BoolCircuit.DAGCircuit.size, BoolCircuit.GatesAcyclic.nil, BoolCircuit.TreeCircuit, BoolCircuit.TreeCircuit.size, BoolCircuit.TreeCircuit.toDAG, BoolCircuit.compileTree_spec}
\difficulty{3}
Let $c$ be a Boolean circuit tree on $n$ variables. The size of its compiled directed
acyclic circuit is at most $n$ plus the number of nodes in $c$.
\end{theorem}

\begin{theorem}[Depth of the circuit DAG translation]
\lean{BoolCircuit.TreeCircuit.toDAG_depth_le}
\label{BoolCircuit.TreeCircuit.toDAG_depth_le}
\leanok
\uses{BoolCircuit.DAGCircuit.depth, BoolCircuit.GatesAcyclic.nil, BoolCircuit.TreeCircuit, BoolCircuit.TreeCircuit.depth, BoolCircuit.TreeCircuit.toDAG, BoolCircuit.compileTree_spec}
\difficulty{1}
For every Boolean circuit $c$ on $n$ variables, the depth of its directed acyclic graph
translation is at most $\mathrm{depth}(c)+1$.
\end{theorem}

\begin{theorem}[Well-formedness of compiled circuits]
\lean{BoolCircuit.TreeCircuit.toDAG_isWellFormed}
\label{BoolCircuit.TreeCircuit.toDAG_isWellFormed}
\leanok
\uses{BoolCircuit.DAGCircuit.IsWellFormed, BoolCircuit.GatesAcyclic.nil, BoolCircuit.TreeCircuit, BoolCircuit.TreeCircuit.toDAG, BoolCircuit.compileTree_spec}
\difficulty{3}
Let $c$ be a Boolean circuit tree. Every gate in its compiled directed acyclic circuit
has no repeated input vertices, and each NOT gate has exactly one input.
\end{theorem}

\begin{theorem}[Fanin of compiled circuits]
\lean{BoolCircuit.TreeCircuit.toDAG_isFaninTwo}
\label{BoolCircuit.TreeCircuit.toDAG_isFaninTwo}
\leanok
\uses{BoolCircuit.DAGCircuit.IsFaninTwo, BoolCircuit.GatesAcyclic.nil, BoolCircuit.TreeCircuit, BoolCircuit.TreeCircuit.maxFanin, BoolCircuit.TreeCircuit.toDAG, BoolCircuit.TreeCircuit.toDAG_isWellFormed, BoolCircuit.compileTree_spec}
\difficulty{4}
Let $c$ be a Boolean circuit tree whose maximum fanin is at most $2$. Then every gate in
its compiled directed acyclic circuit has at most two inputs, no repeated input
vertices, and any NOT gate has exactly one input.
\end{theorem}

\begin{definition}[Unfolding a gate]
\lean{BoolCircuit.unfoldGate}
\label{BoolCircuit.unfoldGate}
\leanok
\uses{BoolCircuit.DAGGate, BoolCircuit.TreeCircuit}
Let $g$ be a gate, let $t$ assign a Boolean circuit tree to each natural-numbered vertex
and polarity, and let $p$ be a polarity. An AND gate unfolds to a node with polarity $p$
whose children are the trees assigned by $t$ to its argument vertices with polarity $p$;
an OR gate unfolds to a node with polarity opposite to $p$ and those same child
assignments. A NOT gate unfolds to the tree assigned to the first argument, or to vertex
$0$ if there is no argument, with polarity opposite to $p$.
\end{definition}

\begin{definition}[Fuel-bounded vertex unfolding]
\lean{BoolCircuit.unfoldVertex}
\label{BoolCircuit.unfoldVertex}
\leanok
\uses{BoolCircuit.DAGGate, BoolCircuit.TreeCircuit, BoolCircuit.constGate, BoolCircuit.unfoldGate}
Given a natural number $n$, a list of gates, a fuel bound, a vertex $v$, and a polarity
$p$, unfold the vertex recursively as follows. With zero fuel, return an empty AND node.
With positive fuel, if $v<n$, return the literal at variable $v$ with polarity $p$;
otherwise unfold the gate at index $v-n$, using an empty AND gate if that index is
outside the list, and recursively unfold its argument vertices with one less unit of
fuel.
\end{definition}

\begin{definition}[Expansion of a circuit into a tree]
\lean{BoolCircuit.DAGCircuit.toTree}
\label{BoolCircuit.DAGCircuit.toTree}
\leanok
\uses{BoolCircuit.DAGCircuit, BoolCircuit.TreeCircuit, BoolCircuit.unfoldVertex}
For a Boolean circuit represented by an acyclic directed graph, its associated tree
circuit is obtained by unfolding the output vertex and recursively copying the gates and
their inputs.
\end{definition}

\begin{theorem}[Depth of a mapped list of circuits]
\lean{BoolCircuit.maxDepth_map}
\label{BoolCircuit.maxDepth_map}
\leanok
\uses{BoolCircuit.TreeCircuit, BoolCircuit.TreeCircuit.depth, BoolCircuit.TreeCircuit.maxDepth, BoolCircuit.TreeCircuit.maxDepth_cons}
\difficulty{4}
Let $n$ be a natural number, let $\ell$ be a list of elements of a type $\alpha$, and
let $f$ assign to each element a circuit on $n$ variables. The maximum depth among the
circuits obtained by applying $f$ to the elements of $\ell$ equals the maximum, with
value $0$ for an empty list, of their individual depths.
\end{theorem}

\begin{theorem}[Size of a mapped list of circuits]
\lean{BoolCircuit.sumSize_map}
\label{BoolCircuit.sumSize_map}
\leanok
\uses{BoolCircuit.TreeCircuit, BoolCircuit.TreeCircuit.size, BoolCircuit.TreeCircuit.sumSize, BoolCircuit.TreeCircuit.sumSize_cons}
\difficulty{4}
Let $n$ be a natural number, let $\alpha$ be a type, let $l$ be a finite list of
elements of $\alpha$, and let $f$ map each element of $\alpha$ to a Boolean circuit on
$n$ variables. The sum of the sizes of the circuits $f(a)$ as $a$ ranges over $l$,
counted with repetitions, equals the sum of the sizes of the mapped list of circuits.
\end{theorem}

\begin{theorem}[Maximum fan-in of a mapped list]
\lean{BoolCircuit.maxFaninL_map}
\label{BoolCircuit.maxFaninL_map}
\leanok
\uses{BoolCircuit.TreeCircuit, BoolCircuit.TreeCircuit.maxFanin, BoolCircuit.TreeCircuit.maxFaninL, BoolCircuit.TreeCircuit.maxFaninL_cons}
\difficulty{4}
Let $n$ be a natural number, let $\ell$ be a list of elements of a type $\alpha$, and
let $f$ assign a circuit on $n$ variables to each element. The maximum fan-in among the
circuits obtained by applying $f$ to the elements of $\ell$ equals the maximum, with
value $0$ for an empty list, of their individual maximum fan-ins.
\end{theorem}

\begin{theorem}[Bound by the maximum of mapped list values]
\lean{BoolCircuit.le_foldr_max}
\label{BoolCircuit.le_foldr_max}
\leanok
\difficulty{4}
Let $l$ be a list of natural numbers, let $f:\mathbb{N}\to\mathbb{N}$, and let $a\in l$.
Then $f(a)$ is at most the maximum of the values of $f$ on the entries of $l$, with the
maximum taken together with $0$.
\end{theorem}

\begin{theorem}[Bound on a sum by a uniform bound]
\lean{BoolCircuit.sum_le_card_mul}
\label{BoolCircuit.sum_le_card_mul}
\leanok
\difficulty{4}
Let $l$ be a finite list of natural numbers, let $f$ map natural numbers to natural
numbers, and let $B$ be a natural number. If $f(a)\le B$ for every $a$ in $l$, then the
sum of the values $f(a)$ over the list, counting repetitions, is at most the length of
$l$ multiplied by $B$.
\end{theorem}

\begin{theorem}[Evaluation of an unfolded gate]
\lean{BoolCircuit.unfoldGate_eval}
\label{BoolCircuit.unfoldGate_eval}
\leanok
\uses{BoolCircuit.DAGGate, BoolCircuit.DAGGate.WellFormed, BoolCircuit.DAGGate.eval, BoolCircuit.TreeCircuit, BoolCircuit.TreeCircuit.eval, BoolCircuit.TreeCircuit.eval_node_false, BoolCircuit.TreeCircuit.eval_node_true, BoolCircuit.unfoldGate}
\difficulty{5}
Let $n$ be a natural number, let $g$ be a well-formed gate, let $t$ assign a Boolean
circuit on $n$ variables to each gate argument and polarity, let $v$ be a list of
Boolean values, and let $x$ be an assignment to the $n$ variables. Suppose that for
every argument $a$ of $g$ and each polarity $q$, the circuit $t(a,q)$ evaluates under
$x$ to the value of $v$ at index $a$ (using false when that index is out of range) if
$q$ is true, and to its negation if $q$ is false. Then unfolding $g$ at polarity $p$
evaluates under $x$ to the value of $g$ on $v$ if $p$ is true, and to its negation if
$p$ is false.
\end{theorem}

\begin{theorem}[Depth bound for gate unfolding]
\lean{BoolCircuit.unfoldGate_depth_le}
\label{BoolCircuit.unfoldGate_depth_le}
\leanok
\uses{BoolCircuit.DAGGate, BoolCircuit.DAGGate.WellFormed, BoolCircuit.TreeCircuit, BoolCircuit.TreeCircuit.depth, BoolCircuit.TreeCircuit.depth_node, BoolCircuit.foldr_max_le, BoolCircuit.maxDepth_map, BoolCircuit.unfoldGate}
\difficulty{4}
Let $n$ be a natural number, let $g$ be a well-formed gate, and let $t$ assign a Boolean
circuit on $n$ variables to each gate argument and polarity. If every circuit $t(a,q)$,
for an argument $a$ of $g$ and polarity $q$, has depth at most $B$, then unfolding $g$
with $t$ at any polarity has depth at most $B+1$.
\end{theorem}

\begin{theorem}[Maximum fan-in bound after unfolding]
\lean{BoolCircuit.unfoldGate_maxFanin_le}
\label{BoolCircuit.unfoldGate_maxFanin_le}
\leanok
\uses{BoolCircuit.DAGGate, BoolCircuit.DAGGate.WellFormed, BoolCircuit.TreeCircuit, BoolCircuit.TreeCircuit.maxFanin, BoolCircuit.TreeCircuit.maxFanin_node, BoolCircuit.foldr_max_le, BoolCircuit.maxFaninL_map, BoolCircuit.unfoldGate}
\difficulty{4}
Let $n$ be a natural number, let $g$ be a well-formed gate, and let $t$ assign a Boolean
circuit on $n$ variables to each gate argument and polarity. If $g$ has at most $k$
arguments and every circuit $t(a,q)$ assigned to an argument of $g$ has maximum fan-in
at most $k$, then the circuit obtained by unfolding $g$ at any polarity has maximum
fan-in at most $k$.
\end{theorem}

\begin{theorem}[Size bound after gate unfolding]
\lean{BoolCircuit.unfoldGate_size_le}
\label{BoolCircuit.unfoldGate_size_le}
\leanok
\uses{BoolCircuit.DAGGate, BoolCircuit.DAGGate.WellFormed, BoolCircuit.TreeCircuit, BoolCircuit.TreeCircuit.size, BoolCircuit.TreeCircuit.size_node, BoolCircuit.sumSize_map, BoolCircuit.sum_le_card_mul, BoolCircuit.unfoldGate}
\difficulty{5}
Let $n$ be a natural number, let $g$ be a well-formed gate, and let $t$ assign a Boolean
circuit on $n$ variables to each gate argument and polarity. Suppose that $g$ has at
most $k$ arguments and every circuit $t(a,q)$ assigned to an argument of $g$ has size at
most $B$. Then, for each polarity $p$, unfolding $g$ using $t$ gives a circuit of size
at most $1+kB$.
\end{theorem}

\begin{theorem}[Defaulted gate lookup]
\lean{BoolCircuit.DAGCircuit.getD_gates}
\label{BoolCircuit.DAGCircuit.getD_gates}
\leanok
\uses{BoolCircuit.constGate}
\difficulty{1}
Let $C$ be a directed acyclic circuit with an ordered list of gates, and let $i$ be a
natural number strictly less than the number of gates. Then the list lookup at index
$i$, using the constant-true gate as a default, equals the gate at index $i$.
\end{theorem}

\begin{theorem}[Arguments precede their vertex]
\lean{BoolCircuit.DAGCircuit.args_lt_vertex}
\label{BoolCircuit.DAGCircuit.args_lt_vertex}
\leanok
\difficulty{3}
Let $C$ be a finite directed acyclic Boolean circuit with $n$ input vertices and a list
of gates in topological order. If $v$ is not an input vertex and $v<n+	ext{number of
gates}$, then every argument of the gate at position $v-n$ has index strictly less than
$v$.
\end{theorem}

\begin{theorem}[Depth decreases along an edge]
\lean{BoolCircuit.DAGCircuit.depthAt_lt_of_arg}
\label{BoolCircuit.DAGCircuit.depthAt_lt_of_arg}
\leanok
\uses{BoolCircuit.DAGCircuit.depthAt, BoolCircuit.DAGCircuit.depthAt_gate, BoolCircuit.DAGCircuit.depths, BoolCircuit.DAGGate.depth, BoolCircuit.le_foldr_max}
\difficulty{4}
Let $C$ be a directed acyclic Boolean circuit on $n$ input vertices. If $v$ is not an
input vertex and $v<n+|C|$, then every vertex $a$ read by the gate at index $v-n$ has
strictly smaller depth than $v$.
\end{theorem}

\begin{theorem}[Evaluation of bounded unfolding]
\lean{BoolCircuit.DAGCircuit.unfoldVertex_eval}
\label{BoolCircuit.DAGCircuit.unfoldVertex_eval}
\leanok
\uses{BoolCircuit.DAGCircuit.IsWellFormed, BoolCircuit.DAGCircuit.args_lt_vertex, BoolCircuit.DAGCircuit.getD_gates, BoolCircuit.DAGCircuit.values, BoolCircuit.DAGCircuit.values_getD_gate, BoolCircuit.DAGCircuit.values_getD_input, BoolCircuit.Lit.eval, BoolCircuit.TreeCircuit.eval, BoolCircuit.unfoldGate_eval, BoolCircuit.unfoldVertex}
\difficulty{5}
Let $C$ be a Boolean directed acyclic circuit whose gates have pairwise distinct input
vertices and whose NOT gates each have exactly one input, and let $x$ be an assignment
of Boolean values to its $n$ input variables. For every natural number $f$, vertex index
$v$, and polarity $p$, if $v<n+	ext{number of gates of }C$ and $v<f$, then unfolding $v$
with fuel $f$ and polarity $p$ evaluates to the value of vertex $v$ under $x$ when $p$
is true, and to its negation when $p$ is false.
\end{theorem}

\begin{theorem}[Evaluation of the unfolded circuit]
\lean{BoolCircuit.DAGCircuit.toTree_eval}
\label{BoolCircuit.DAGCircuit.toTree_eval}
\leanok
\uses{BoolCircuit.DAGCircuit.IsWellFormed, BoolCircuit.DAGCircuit.eval, BoolCircuit.DAGCircuit.toTree, BoolCircuit.DAGCircuit.unfoldVertex_eval, BoolCircuit.TreeCircuit.eval}
\difficulty{3}
Let $C$ be a Boolean directed acyclic circuit whose gates have pairwise distinct input
vertices and whose NOT gates each have exactly one input. For every Boolean assignment
$x$ to its $n$ input variables, the tree obtained by unfolding $C$ computes the same
value as $C$ under $x$.
\end{theorem}

\begin{theorem}[Depth bound for bounded unfolding]
\lean{BoolCircuit.DAGCircuit.unfoldVertex_depth_le}
\label{BoolCircuit.DAGCircuit.unfoldVertex_depth_le}
\leanok
\uses{BoolCircuit.DAGCircuit.IsWellFormed, BoolCircuit.DAGCircuit.args_lt_vertex, BoolCircuit.DAGCircuit.depthAt, BoolCircuit.DAGCircuit.depthAt_gate, BoolCircuit.DAGCircuit.depthAt_lt_of_arg, BoolCircuit.DAGCircuit.getD_gates, BoolCircuit.DAGGate.depth, BoolCircuit.TreeCircuit.depth, BoolCircuit.unfoldGate_depth_le, BoolCircuit.unfoldVertex}
\difficulty{5}
Let $C$ be a Boolean directed acyclic circuit whose gates have pairwise distinct input
vertices and whose NOT gates each have exactly one input. For every natural number $f$,
vertex index $v$, and polarity $p$, if $v<n+	ext{number of gates of }C$ and $v<f$, then
the depth of the tree obtained by unfolding $v$ with fuel $f$ and polarity $p$ is at
most the depth assigned to vertex $v$ by $C$.
\end{theorem}

\begin{theorem}[Depth bound for unfolding a DAG]
\lean{BoolCircuit.DAGCircuit.toTree_depth_le}
\label{BoolCircuit.DAGCircuit.toTree_depth_le}
\leanok
\uses{BoolCircuit.DAGCircuit.IsWellFormed, BoolCircuit.DAGCircuit.depth, BoolCircuit.DAGCircuit.toTree, BoolCircuit.DAGCircuit.unfoldVertex_depth_le, BoolCircuit.TreeCircuit.depth}
\difficulty{1}
Let $C$ be a Boolean DAG circuit whose gates have distinct arguments and whose NOT gates
each have exactly one argument. The depth of the tree obtained by unfolding $C$ is at
most the depth of $C$.
\end{theorem}

\begin{theorem}[Fanin bound for unfolded circuits]
\lean{BoolCircuit.DAGCircuit.unfoldVertex_maxFanin_le}
\label{BoolCircuit.DAGCircuit.unfoldVertex_maxFanin_le}
\leanok
\uses{BoolCircuit.DAGCircuit.IsWellFormed, BoolCircuit.DAGCircuit.args_lt_vertex, BoolCircuit.DAGCircuit.getD_gates, BoolCircuit.TreeCircuit.maxFanin, BoolCircuit.unfoldGate_maxFanin_le, BoolCircuit.unfoldVertex}
\difficulty{4}
Let $C$ be a Boolean directed acyclic circuit with $n$ inputs, whose gates have distinct
argument vertices and whose NOT gates each have exactly one argument. Suppose every gate
has at most $k$ arguments. For every fuel bound, polarity, and vertex
$v<n+|C.	ext{gates}|$, the maximum fanin of the tree obtained by unfolding $v$ is at
most $k$.
\end{theorem}

\begin{theorem}[Fanin of the unfolded circuit]
\lean{BoolCircuit.DAGCircuit.toTree_maxFanin_le}
\label{BoolCircuit.DAGCircuit.toTree_maxFanin_le}
\leanok
\uses{BoolCircuit.DAGCircuit.IsWellFormed, BoolCircuit.DAGCircuit.toTree, BoolCircuit.DAGCircuit.unfoldVertex_maxFanin_le, BoolCircuit.TreeCircuit.maxFanin}
\difficulty{1}
Let $C$ be a Boolean directed acyclic circuit whose gates have pairwise distinct input
vertices and whose NOT gates each have exactly one input. If every gate has at most $k$
input vertices, then the maximum fanin of the tree obtained by unfolding $C$ is at most
$k$.
\end{theorem}

\begin{theorem}[Size bound for unfolded circuits]
\lean{BoolCircuit.DAGCircuit.unfoldVertex_size_le}
\label{BoolCircuit.DAGCircuit.unfoldVertex_size_le}
\leanok
\uses{BoolCircuit.DAGCircuit.IsWellFormed, BoolCircuit.DAGCircuit.args_lt_vertex, BoolCircuit.DAGCircuit.depthAt, BoolCircuit.DAGCircuit.depthAt_gate, BoolCircuit.DAGCircuit.depthAt_lt_of_arg, BoolCircuit.DAGCircuit.getD_gates, BoolCircuit.DAGGate.depth, BoolCircuit.TreeCircuit.size, BoolCircuit.unfoldGate_size_le, BoolCircuit.unfoldVertex}
\difficulty{5}
Let $C$ be a Boolean directed acyclic circuit with $n$ inputs, whose gates have distinct
argument vertices and whose NOT gates each have exactly one argument. Suppose every gate
has at most $k$ arguments. For every fuel bound, polarity, and vertex
$v<n+|C.	ext{gates}|$ with $v$ less than the fuel bound, the unfolded tree at $v$ has
size at most $(k+1)^{d(v)}$, where $d(v)$ is the depth assigned to vertex $v$ in $C$.
\end{theorem}

\begin{theorem}[Size of the unfolded circuit]
\lean{BoolCircuit.DAGCircuit.toTree_size_le}
\label{BoolCircuit.DAGCircuit.toTree_size_le}
\leanok
\uses{BoolCircuit.DAGCircuit.IsWellFormed, BoolCircuit.DAGCircuit.depth, BoolCircuit.DAGCircuit.toTree, BoolCircuit.DAGCircuit.unfoldVertex_size_le, BoolCircuit.TreeCircuit.size}
\difficulty{1}
Let $C$ be a Boolean directed acyclic circuit whose gates have pairwise distinct input
vertices and whose NOT gates each have exactly one input. If every gate has at most $k$
input vertices, then the size of the tree obtained by unfolding $C$ is at most
$(k+1)^{	ext{depth}(C)}$.
\end{theorem}
